\scanpage{14}
\mspage{7})

qu'on compose encore comme précédemment. On a bien \deleted{\unsure{deux|un}} \unsure{foncteurs|foncteur}
\diagram{\begin{tikzcd}
\deleted{\unsure{C\longrightarrow}}&\widetilde C\arrow[r]&W^{-1}C
\end{tikzcd}}

\deleted{Le premier de ces deux foncteurs transforme bien les éléments de $W$ en flèches inversibles. \deleted{On a} Pour
\diagram{\begin{tikzcd}
x\arrow[r,"f","\sim"']&y
\end{tikzcd}}
\unsure{$f\in W$}, on considère
\[\text{soit }F=[f]\]
\diagram{\begin{tikzcd}
y\arrow[dr,"1_y"']&x\arrow[d,"f","\sim"']\\
&y
\end{tikzcd}}
\[\text{soit }G=[1_y,f].\]
\unsure{Et on trouve les composés} \unsure{et on a donc}
\[\uline{GF}=\unsure{[\,]^{x}}\]
}
\unsure{définissons le} $=$
\diagram{\begin{tikzcd}C\arrow[r]&\widetilde C\end{tikzcd}}
qui va être un peu plus délicat. Soit $f:x\longrightarrow y$. On choisit une factorisation
\diagram{\begin{tikzcd}[column sep=large,row sep=large]
x\arrow[dr,hook,"f_0"']&y\arrow[dr,bend left=30,dotted,"i"]\arrow[dr,dashed,"\beta","\sim"']&\\
&\bar y\arrow[u,"\unsure{\alpha}","\sim"']\arrow[r,dotted,"\unsure{g}","\sim"']&\bar y'
\end{tikzcd}}
\inserted{\unsure{donc $\bar y$ cof}}
et, avec le \unsure{w.k lifting lemma} (ici $y$ étant cofibrant) on peut compléter par des pointillés, où $\beta$ a une section $i$.

On considère alors $[g f_0,i]$, il faut montrer que ça ne dépend pas des choix faits. Pour un choix $[f',i']$
\diagram{\begin{tikzcd}
x\arrow[r,dashed,"f"]\arrow[dr,"f'"']&y\arrow[d,"i'","\sim"',shift right=1]\\
&y'\arrow[u,dashed,"\alpha","\sim"',shift right=1]
\end{tikzcd}}
tel qu'il existe \deleted{\unsure{un|une}} une rétraction $\alpha$ de $\widetilde y$ sur $y$, telle que
\[\alpha f'=f.\]

\margin{On s'est ramené aux factorisations $f=\unsure{\alpha}f_0$, $\unsure{f_0}\in\unsure{\operatorname{cof}}$ et $\unsure{\alpha}\in W$, telles que $\unsure{\alpha}$ ait une section.}

Considérons deux tels choix, il faut voir qu'ils sont dans une $=$ composante connexe de $\underline{\operatorname{Ch}}_0(x,y)$. Mais pour deux factorisations d'une flèche $f:x\longrightarrow y$ en $\alpha f_0$ et $\alpha'f'_0$, $f_0,f'_0\in\operatorname{Cof}$, $\alpha,\alpha'\in W$, \deleted{\unsure{$x\longrightarrow y$}} on peut les coiffer par une troisième \deleted{\unsure{$f''$}} $\alpha'',f''_0$ qui \deleted{\unsure{\diagram{\begin{tikzcd}x\arrow[r]&y''\end{tikzcd}}}} \deleted{\unsure{permet de}} \inserted{\unsure{reçoit|réduit}} les deux \unsure{comme on voit dans les diagrammes}

\margin{Ce foncteur est bijectif sur les objets, surjectif sur les flèches, \unsure{comme préc.}. On vérifie aisément [il suffit pour cela \unsure{dans un produit} de voir que pour tt $f\in W$, $f:x\longrightarrow y$, $[f]^{-1}$ provient de $\widetilde C$, on prend
\diagram{\begin{tikzcd}
x\arrow[r,"\unsure{f}"]\arrow[dr,"\sim"']&y\arrow[d,shift left=1,"\alpha"]\\
&\unsure{\bar y}\arrow[u,very thick,shift left=1,"\unsure{i}","\sim"']
\end{tikzcd}}
\unsure{(les deux isos)}]
\diagram{\begin{tikzcd}
y\arrow[dr,"\unsure{1_y}"']&\unsure{x}\arrow[d,"\unsure{f}","\sim"',shift right=1]\\
&y\arrow[u,"\unsure{\alpha}"',shift right=1]
\end{tikzcd}}
}

\figure{\begin{tikzpicture}[x=.4cm,y=.4cm]
\draw (0,.2) .. controls (-1.4,.4) and (-.8,3.2) .. (1.8,5.5) .. controls (3.8,7.3) and (5.4,7.8) .. (5.1,5.4) .. controls (4.8,2.1) and (1.8,-.5) .. (0,.2);
\end{tikzpicture}}
\mspage{13}
